\chapter{结论与展望}

\section{研究结论}

本可行性研究报告对“家庭储能优化调度系统”进行了全方位的深度论证。基于详尽的市场调研、严谨的数学建模和多场景的仿真实验，得出以下核心结论：

\begin{enumerate}
      \item \textbf{技术架构先进且可控}：
            项目采用的“端-云-边”协同架构（Flutter + Flatten + GCP）经论证具备极高的可扩展性。通过引入混合整数规划 (MIP) 模型，系统能够在满足电池 SOC 约束和防逆流约束的前提下，实现毫秒级的最优策略求解。网络拓扑分析显示，基于私有 VPC 的部署方案能有效隔离外部攻击，技术风险处在可控范围内。

      \item \textbf{经济效益具有吸引力}：
            敏感性分析表明，在当前电池成本（\$500/kWh）和加州分时电价下，系统的投资回报率 (ROI) 已达 \textbf{16\%}，静态回收期约为 6.25 年。随着未来电池成本下降至 \$300/kWh，ROI 有望突破 28\%。与传统规则控制策略相比，本系统的 AI 优化算法能额外提升 15\%-20\% 的节能收益。

      \item \textbf{运维与社会价值显著}：
            CI/CD 流水线的引入确保了软件的高质量快速交付，用户泳道图分析证实了产品操作的便捷性。在社会层面，单户年均 1.8 吨的碳减排潜力使其完全契合全球“碳中和”战略，具备接入虚拟电厂 (VPP) 参与电网辅助服务的巨大潜力。
\end{enumerate}

综上所述，本项目在技术上成熟、经济上合理、法律上合规，具有极高的落地可行性和广阔的商业前景。

\section{未来展望}

尽管当前方案已具备极高的可行性，但面向未来，系统仍有演进空间：
\begin{enumerate}
      \item \textbf{引入 V2G (Vehicle-to-Grid) 技术}：随着电动汽车普及，将电动汽车电池纳入家庭储能资源池，可将可用调度容量提升 5-10 倍。
      \item \textbf{强化边缘计算能力}：在本地网关部署轻量级 TFLite 模型，实现断网情况下的离线预测与调度，进一步提升系统韧性。
      \item \textbf{基于区块链的 P2P 能源交易}：利用区块链技术，实现社区内邻里间的分布式能源直接交易，打破传统售电模式。
\end{enumerate}

建议项目组立即启动第一阶段的详细设计工作，争取在 3 个月内发布 MVP 版本。
